\documentclass[11pt]{article}
\usepackage[utf8]{inputenc}
\usepackage{amsmath,amssymb}
\usepackage[margin=1in]{geometry}
\usepackage{hyperref}
\emergencystretch=3em

\title{Headline statements, part IV: leading rates of the chart posterior ratio}
\author{Claude, with Daniel Murfet}
\date{2026-09-07}

\begin{document}
\maketitle
\sloppy

\begin{abstract}
Paper-facing wrappers for units~166--168 (file \texttt{HeadlinePosteriorRates.lean}):
the quotient calculus for power--log leading terms, the two-term expansion with power-saving
remainder at the leading exponent, the constant posterior limit $Z[\varphi]/Z[1]\to\eta_\varphi(0)/\eta_1(0)$
for an observable not vanishing at the corner, and the $1/\log N$ decay
$(Z[\varphi]/Z[1])\log N\to B^\varphi_p/A^1_p$ for one that does. The mirror's remark on
cor:empirical\_expectation was extended accordingly (staging copy committed; Overleaf sync pending
authorisation). Zero \texttt{sorry}/\texttt{axiom}.
\end{abstract}

\section{Scope}
Deterministic, one chart, $d=2$, equal starting exponents, fixed phase data, amplitudes with
$\eta_1(0)>0$. Not claimed: a power improvement from corner vanishing (false in general), the next
exponent under coordinatewise divisibility, unequal starting exponents, the stochastic $1/\log N$
regime.

\section{Lean notes}
Wrappers only; compiled on the first attempt.

\end{document}
